\documentclass[10pt,a4paper]{amsart}
\usepackage[margin=19mm]{geometry}
\usepackage{amsmath,amssymb,mathtools}
\usepackage[scr=rsfs,cal=euler]{mathalfa}
\usepackage{xfrac}
\usepackage[unicode,hidelinks,bookmarks=false]{hyperref}
\newcommand{\ie}{i.e.\ }
\newcommand{\cf}{cf.\ }
\newcommand{\resp}{respectively\ }
\newcommand{\Subsection}{\subsection}
\providecommand{\nobreakdash}{\nobreak\hbox{-}}
\setlength{\emergencystretch}{2em}
\multlinegap0pt
\usepackage{amssymb}
\usepackage{mathtools}
\usepackage{xcolor}
\usepackage[shortlabels]{enumitem}
\usepackage{dsfont}
\newtheorem{corollary}{Corollary}
\newtheorem{theorem}{Theorem}
\newtheorem{prop}{Proposition}
\newtheorem{lemma}{Lemma}
\newcommand{\G}{\mathnormal{\mathbb{T}^{m}\times\mathbb{R}^n}}
\newcommand{\Imag}{\mathrm{Im}}
\newcommand{\N}{\mathbb{N}}
\newcommand{\R}{\mathbb{R}}
\newcommand{\Real}{\mathrm{Re}}
\renewcommand{\S}{\mathcal{S}_{\sigma,\mu}(\G)}
\newcommand{\Smu}{\mathscr{F}_{\mu}(\T^1\times\R)}
\newcommand{\T}{\mathbb{T}}
\newcommand{\Z}{\mathbb{Z}}
\newcommand{\g}{\mathbb{T}^1\times\mathbb{R}}
\newcommand{\s}{\mathcal{S}(\G)}
\newcommand{\sigmac}{\sigma_{\scalebox{0.5}{$0$}}}
\title{Global hypoellipticity on time-periodic Gelfand-Shilov spaces via non-discrete Fourier analysis}
\author{Andr\'e Pedroso Kowacs, Pedro Meyer Tokoro}
\date{Source: arXiv:2506.13475v1 (CC BY 4.0)}
\begin{document}
\maketitle

\noindent\emph{Source and licence.} Global hypoellipticity on time-periodic Gelfand-Shilov spaces via non-discrete Fourier analysis. Version arXiv:2506.13475v1 (CC BY 4.0), distributed by arXiv under the Creative Commons Attribution 4.0 International licence, http://creativecommons.org/licenses/by/4.0/, as recorded in the arXiv licence metadata for this exact version and shown on the abstract page https://arxiv.org/abs/2506.13475v1. Full licence notice: \texttt{licenses/arxiv-2506.13475-CC-BY-4.0.txt}.\par\smallskip

\noindent\textit{Editorial scope.} Counted unit: the article's prop prop (source0.tex lines 1287--1289). The article's complete original proof of it is delivered with it (source0.tex lines 1290--1395, one \texttt{proof} environment of 106 lines). No mathematics is written, repaired or re-derived: the statement and the proof are the article's own lines, in the article's own notation, and no retained line is substituted or re-flowed.  Retained uncounted as the source-local context the proof needs: lines 495--497; lines 671--674; lines 961--969; lines 1448--1453; lines 1459--1464; lines 1470--1475; lines 1502--1522.  Nothing was omitted from any retained span, so the package carries no omission record.  No retained line contains a citation, so the wrapper carries no bibliography and no bibliography entry.  No retained line contains a figure environment, an image-inclusion command or drawing code, so no image is delivered and the package declares no asset dependency.  Editorial formatting only, in addition: an AMS \texttt{amsart} preamble that restates the article's own declarations and macros used by the retained lines, namely the environments \texttt{corollary}, \texttt{theorem}, \texttt{prop}, \texttt{lemma} on separate counters, together with the standard packages those lines need.  The retained environments print in this document as Corollary 1, Theorem 1, Proposition 1, Lemma 1 in order of appearance, whereas the article prints the same items with the numbering of its own full document.  The complete native source of the article as distributed in the arXiv e-print for this exact version is bundled unchanged as \texttt{sources/179948/source0.tex}.
    \begin{corollary}\label{coro_gelf_transf} 
        Let $f\in L^1(\T^m\times\R^n)$. Then $f\in\S$ if and only if $\mathcal{F}_{\R^n}f\in\S$.
    \end{corollary}

\begin{theorem}\label{theo_gh_cte_Nleq1}
    Let $P$ be a differential operator on $\T^1\times\R$ with constant coefficients and no mixed terms, such that  $$P=p\left(D_x\right)+q\left(D_t\right),$$
   where $p$ and $q$ are complex polynomials and $\deg(p)\leq 1$. Then $P$ is $\mathcal{S}_{\sigma,\mu}$-globally hypoelliptic if and only if its symbol does not vanish on $\Z\times \R$.
\end{theorem}

\begin{corollary}\label{coro-tube-cte}
    The operator
    \[P=\partial_t + (a_0+ib_0)\partial_x + q_0,\]
    where $q_0\in\mathbb{C}$, $a_0,b_0\in\R$ is $\mathscr{F}_\mu$-globally hypoelliptic if and only if
    \begin{enumerate}
        \item[(i)] $b_0\neq 0$ and $\frac{a_0}{b_0}\Real(q_0)+\Imag(q_0)\notin\mathbb{Z}$;
        \item[(ii)] $b_0=0$ and either $\Real(q_0)\neq 0$, or both $a_0=\Real(q_0)=0$ and $\operatorname{Im}(q_0)\notin\Z$.
    \end{enumerate}
\end{corollary}

\begin{lemma}\label{derivative-reciprocal}
    Let $g\in C^{\infty}(\R)$ be such that $g(t)\neq 0$ for every $t\in\R$. Then for every $n\in\N$ we have that
    \begin{equation*}
        \partial_t^n [g(t)^{-1}]=\sum_{k=1}^n(-1)^k\binom{n+1}{k+1}\frac{1}{g(t)^{k+1}}\partial_t^n[g(t)^k],\quad \forall t\in\R.
    \end{equation*}
\end{lemma}

\begin{lemma}\label{lemma_algebraic_1}
    If $\tau=(\tau_1,\dots,\tau_N)\in\Delta(N)$ then, for any $\sigma\geq 1$, we have that
    \begin{equation*}
        |\tau|!^{\sigma}\prod_{\ell=1}^{N}\ell!^{(\sigma-1)\tau_{\ell}}\leq |\tau|!N!^{\sigma-1}
    \end{equation*}
\end{lemma}

\begin{lemma}\label{lemma_algebraic_2}
    For each $\N\in\N$, $\tau\in\N_0^N$ and $R\in\R$ we have that
    \begin{equation*}
        \sum_{\Delta(N)}\frac{|\tau|!}{\tau!}R^{|\tau|}=R(1+R)^{N-1}.
    \end{equation*}
\end{lemma}

\begin{lemma}\label{lemmaodesol}
	Let $g,\theta \in C^{\infty}(\T^1)$, and set $\theta_0 = \frac{1}{2\pi}\int_0^{2\pi}\theta(t)\,\mathrm{d}t$. If $\theta_0\not\in i\mathbb{Z}$, then the differential equation
	\begin{equation}\label{ode}
		\partial_t u(t)+\theta(t)u(t)=g(t),\,\qquad t\in\T^1,
	\end{equation}
	admits a unique solution in $C^\infty(\T^1)$ given by
	\begin{equation}\label{sol-}
		u(t) = \frac{1}{1-e^{-2\pi\theta_0}}\int_0^{2\pi}g(t-s)e^{-\int_{t-s}^t\theta(\tau)\,\mathrm{d}\tau}\,\mathrm{d}s,
	\end{equation} 
	or equivalently by
	\begin{equation}\label{sol+}
		u(t) = \frac{1}{e^{2\pi\theta_0}-1}\int_0^{2\pi}g(t+s)e^{\int_{t}^{t+s}\theta(\tau)\,\mathrm{d}\tau}\,\mathrm{d}s.
	\end{equation}
	
	If $\theta_0\in i\mathbb{Z}$, then equation \eqref{ode} admits infinitely many solutions given by
	\begin{equation*}
		u_\lambda(t) = \lambda e^{-\int_0^t\theta(\tau)\,\mathrm{d}\tau}+\int_0^t g(s)e^{-\int_s^t\theta(\tau)\,\mathrm{d}\tau}\,\mathrm{d}s,
	\end{equation*}
	for every $\lambda\in\mathbb{R}$, if and only if
	$$ \int_0^{2\pi}g(t)e^{\int_0^t\theta(\tau)\,\mathrm{d}\tau}\,\mathrm{d}t=0.$$
\end{lemma}

\begin{prop}
    Suppose that $\tilde P$ is $\mathscr{F}_{\mu}$-globally hypoelliptic, $b\not\equiv 0$ and $b$ does not change sign. Then, $P$ is $\mathscr{F}_{\mu}$-globally hypoelliptic for any $\mu\geq 1/2$.
\end{prop}

\begin{proof}
    Suppose that $Pu=f\in\Smu$. First note that due to the inclusions $\mathscr{G}^{\sigma_1}(\T^1)\subset \mathscr{G}^{\sigma_2}(\T^1)$ and $\mathcal{S}_{\sigma_1,\mu}(\g)\subset\mathcal{S}_{\sigma_2,\mu}(\g)$ for $\sigma_1\leq \sigma_2$, we can assume without loss of generality that $\sigmac=\sigma$, where  $f\in\s$. Then, by taking the partial Fourier transform with respect to the second variable, we see that $\widehat{u}(t,\xi)$ must be the only solution to $\widehat{Pu}(t,\xi) = \widehat{f}(t,\xi)$ given by any of the two equivalent formulas in Lemma \ref{lemmaodesol}. Suppose first $b(t)\geq 0$, which implies $ b_0>0$. For $\xi\geq 0$ consider the solution $\widehat{u}(t,\xi)$ given by the formula
    \begin{align*}
        \widehat{u}(t,\xi)=\frac{1}{e^{2\pi i(\xi c_0-iq_0)}-1}\int_{0}^{2\pi}\widehat{f}(t+s,\xi)e^{qs}e^{i\xi a_0s-\xi\int_{t}^{t+s}b(w)\,\mathrm{d}w}\,\mathrm{d}s,
    \end{align*}
    and for $\xi<0$, consider the equivalent formula
        \begin{align*}
        \widehat{u}(t,\xi)=\frac{1}{1-e^{-2\pi i(\xi c_0-iq_0)}}\int_{0}^{2\pi}\widehat{f}(t-s,\xi)e^{-qs}e^{-i\xi a_0s+\xi\int_{t-s}^{t}b(w)\,\mathrm{d}w}\,\mathrm{d}s.
    \end{align*}

    Set 
    \begin{equation*}
        \mathcal{H}_1(\xi)=1-e^{-2\pi i (\xi c_0-i q_0)}\quad\text{and}\quad 
         \mathcal{H}_2(t,s,\xi)=e^{-qs+\xi\left(\int_{t-s}^t b(w)\,\mathrm{d}w-ia_0s\right)},
    \end{equation*}
    for every $t,s\in [0,2\pi]$, $\xi\in\R$.
    Notice that $\mathcal{H}_1$ is analytic and therefore belongs to $\mathscr{G}^\sigma(\R)$.

    Since $\tilde P$ is $\mathscr{F}_{\mu}$-globally hypoelliptic, by Corollary \ref{coro-tube-cte} we have that $|k+c_0\xi-iq_0|\neq 0$ for every $(k,\xi)\in\Z\times \R$. From the proof of Theorem \ref{theo_gh_cte_Nleq1} we conclude that there exists $\delta>0$ such that 
    $|k+c_0\xi-iq_0|\geq \delta$, for every $(k,\xi)\in\Z\times \R$.

Note that by continuity that this implies that there exists $\varepsilon>0$ such that $|\mathcal{H}_1(\xi)|\geq\varepsilon$, for every $\xi\in\R$. Also, by Lemma \ref{derivative-reciprocal}, for $\ell\geq 1$ we have that
    \begin{align*}
        \partial_\xi^{\ell}[\mathcal{H}_1(\xi)^{-1}]&=\sum_{k=1}^\ell(-1)^k\binom{\ell+1}{k+1}\frac{1}{\mathcal{H}_1(\xi)^{k+1}}\partial_\xi^\ell[\mathcal{H}_1(\xi)^k]\\
        &=\sum_{k=1}^\ell(-1)^k\binom{\ell+1}{k+1}\frac{1}{\mathcal{H}_1(\xi)^{k+1}}\partial_\xi^\ell\left[\sum_{j=0}^k\binom{k}{j}(-1)^{k-j}e^{-2\pi ij(\xi c_0-iq_0)}\right]\\
        &=\sum_{k=1}^\ell(-1)^k\binom{\ell+1}{k+1}\frac{1}{\mathcal{H}_1(\xi)^{k+1}}\sum_{j=0}^k\binom{k}{j}(-1)^{k-j+\ell}(2\pi i j c_0)^\ell e^{-2\pi ij(\xi c_0-iq_0)},
    \end{align*}
for every $\xi\in\R$. Hence 
        \begin{align}
        \left|\partial_\xi^{\ell}[\mathcal{H}_1(\xi)^{-1}]\right| & \leq(2\pi \ell |c_0|)^\ell\sum_{k=1}^\ell\binom{\ell+1}{k+1}\varepsilon^{-(k+1)}\sum_{j=0}^k\binom{k}{j} e^{-2\pi j \Real(q_0)} \notag \\
        &\leq (2\pi |c_0|)^\ell e^\ell \ell!\sum_{k=1}^\ell\binom{\ell+1}{k+1}\varepsilon^{-(k+1)}\left(1+e^{-2\pi  \Real(q_0)}\right)^k \notag \\
        &\leq \varepsilon^{-1}(2\pi |c_0|)^\ell e^\ell \ell!\sum_{k=0}^{\ell}\binom{\ell}{k}\frac{\ell+1}{k+1}\left(\frac{1+e^{-2\pi  \Real(q_0)}}{\varepsilon}\right)^k \notag \\
        & \leq \varepsilon^{-1}(4\pi |c_0|)^\ell e^\ell \ell!(\ell+1)\left(1+\frac{1+e^{-2\pi  \Real(q_0)}}{\varepsilon}\right)^\ell, \label{H1}
    \end{align}
    for every $\xi<0$ and $\ell\in\N_0$, since $e^{2\pi b_0\xi}\leq 1$ and where we used
    \[\sum_{k=1}^{\ell}\binom{\ell+1}{k+1}\leq \sum_{k=0}^{\ell}\binom{\ell}{k}\frac{\ell+1}{k+1}\leq (\ell+1)\sum_{k=0}^{\ell}\binom{\ell}{k} = (\ell+1)2^\ell.\]
    
    Also, note that since $b(t)\geq 0$ for every $t\in\T^1$, we have that $\int_{t-s}^tb(w)\,\mathrm{d}w\geq 0$ and therefore
    \begin{equation*}
        |\mathcal{H}_2(t,s,\xi)|\leq e^{2\pi|\Real(q_0)|},
    \end{equation*}
for every $t,s\in [0,2\pi]$, $\xi<0$.
    
By the Fa\`a di Bruno's Formula, we have that
    \begin{align}\label{H2_1}
    \partial_t^{\ell_1}\mathcal{H}_2(t,s,\xi) = \mathcal{H}_2(t,s,\xi)\sum_{\tau\in\Delta(\ell_1)}\frac{\ell_1!}{\tau!}\xi^{|\tau|}\prod_{j_1=1}^{\ell_1}\left(\frac{1}{j_1!}\partial_t^{j_1-1}[b(t)-b(t-s)]\right)^{\tau_{j_1}}. 
    \end{align}
    
    Moreover, it is easy to check that 
    \begin{align}\label{H2_2}
        \partial_\xi^{\ell_2'}\mathcal{H}_2(t,s,\xi) = \mathcal{H}_2(t,s,\xi)\left(\int_{t-s}^tb(w)\,\mathrm{d}w-ia_0s\right)^{\ell_2'}.
    \end{align}
    
    Hence, combining \eqref{H2_1} and \eqref{H2_2}, we obtain by the Leibniz's rule that
    \begin{align*}
    \partial_\xi^{\ell_2'}\partial_t^{\ell_1}\mathcal{H}_2(t,s,\xi) =&\sum_{\tau\in\Delta(\ell_1)}\sum_{j_2=0}^{\min\{\ell_2',|\tau|\}}\binom{\ell_2'}{j_2}\mathcal{H}_2(t,s,\xi)\left(\int_{t-s}^tb(w)\,\mathrm{d}w-ia_0s\right)^{\ell_2'-j_2}\\&\times\frac{|\tau|!}{(|\tau|-j_2)!}\xi^{|\tau|-j_2}
    \frac{\ell_1!}{\tau!}\prod_{j_1=1}^{\ell_1}\left(\frac{1}{j_1!}\partial_t^{j_1-1}[b(t)-b(t-s)]\right)^{\tau_{j_1}}. 
\end{align*}

But note that
\begin{equation*}
    \left|\frac{1}{j_1!}\partial_t^{j_1-1}[b(t)-b(t-s)]\right|\leq \frac{1}{j_1!}2C_{0,b}C_{1,b}^{j_1-1}(j_1-1)!^\sigma\leq 2C_{0,b}C_{1,b}^{j_1-1}j_1!^{\sigma-1}, 
\end{equation*}
for some $C_{0,b},C_{1,b}\geq 1$. Hence
\begin{equation*}
   \prod_{j_1=1}^{\ell_1}\left(\frac{1}{j_1!}\partial_t^{j_1-1}[b(t)-b(t-s)]\right)^{\tau_{j_1}}\leq (2C_{0,b})^{|\tau|}C_{1,b}^{\ell_1}\prod_{j_1=1}^{\ell_1}j_1!^{(\sigma-1)\tau_{j_1}}. 
\end{equation*}

Therefore we can estimate
\begin{align}
    \left|\partial_\xi^{\ell_2'}\partial_t^{\ell_1}\mathcal{H}_2(t,s,\xi)\right| \leq &\   |\mathcal{H}_2(t,s,\xi)|\sum_{\tau\in\Delta(\ell_1)}\sum_{j_2=0}^{\min\{\ell_2',|\tau|\}}\binom{\ell_2'}{j_2}|s|^{\ell_2'-j_2}C_{b,a_0}^{\ell_2'-j_2}\frac{|\tau|!}{(|\tau|-j_2)!}\notag\\
    &\  \times |\xi|^{|\tau|-j_2}
    \frac{\ell_1!}{\tau!}(2C_{0,b})^{|\tau|}C_{1,b}^{\ell_1}\prod_{j_1=1}^{\ell_1}j_1!^{(\sigma-1)\tau_{j_1}}\notag\\
    \leq &\  e^{2\pi|\Real(q)|}\sum_{j_2=0}^{\min\{\ell_2',|\tau|\}}\binom{\ell_2'}{j_2}(2\pi)^{\ell_2'}C_{b,a_0}^{\ell_2'}|\xi|^{\ell_1}C_{1,b}^{\ell_1}\notag\\
    &\ \times\sum_{\tau\in\Delta(\ell_1)} 
    \frac{\ell_1!}{\tau!}(2C_{0,b})^{|\tau|}|\tau|!^\sigma\prod_{j_1=1}^{\ell_1}j_1!^{(\sigma-1)\tau_{j_1}}\notag\\
        \leq &\ e^{2\pi|\Real(q)|}(4\pi C_{b,a_0})^{\ell_2'}|\xi|^{\ell_1}C_{1,b}^{\ell_1}\sum_{\tau\in\Delta(\ell_1)} 
    \frac{\ell_1!}{\tau!}(2C_{0,b})^{|\tau|}|\tau|!\ell_1!^{\sigma-1}\notag\\
    \leq &\ e^{2\pi|\Real(q)|}(4\pi C_{b,a_0})^{\ell_2'}|\xi|^{\ell_1}C_{1,b}^{\ell_1}\ell_1!^\sigma 2C_{0,b}(1+2C_{0,b})^{\ell_1-1},\label{H2}
\end{align}
for every $t,s\in[0,2\pi]$ and $\xi<0$, where $C_{b,a_0} = \displaystyle\max_{w\in[0,2\pi]} |b(w)| + |a_0|$ and in the last two lines we applied Lemmas \ref{lemma_algebraic_1} and \ref{lemma_algebraic_2}.

Next note that
\begin{align}\label{eq-util-final-proof}
    \partial_\xi^\beta\partial_t^{\gamma}\widehat{u}(t,\xi)=\sum_{|\ell'|=\beta} \binom{\beta}{\ell'}\partial_\xi^{\ell'_1}[\mathcal{H}_1(\xi)^{-1}]\sum_{|\ell|=\gamma}\binom{\gamma}{\ell}\int_0^{2\pi}\partial_\xi^{\ell'_2}\partial_t^{\ell_1}\mathcal{H}_2(t,s,\xi)\partial_\xi^{\ell'_3}\partial_t^{\ell_2}\widehat{f}(t-s,\xi)\,\mathrm{d}s,
\end{align}
for every $(t,\xi)\in\T^1\times\R$. Since $\widehat{f}\in\s$, there exist $L>0$, $C_{0,f},C_{1,f}\geq 1$ such that
\begin{equation*}
    e^{L|\xi|^{\frac{1}{\mu}}}|\partial_\xi^\beta\partial_t^\gamma \widehat{f}(t,\xi)|\leq C_{0,f}C_{1,f}^{\beta+\gamma}\beta!^\mu\gamma!^\sigma,
\end{equation*}
for every $(t,\xi)\in\T^1\times\R$, $\beta,\gamma\in\N_0$.
Applying \eqref{H1}, \eqref{H2} and the inequality above to $\eqref{eq-util-final-proof}$, yields
\begin{align*}
    &e^{\frac{L}{2}|\xi|^{\frac{1}{\mu}}}|\partial_\xi^\beta\partial_t^{\gamma}\widehat{u}(t,\xi)|\\
    &\leq \sum_{|\ell'|=\beta} \binom{\beta}{\ell'}\varepsilon^{-1}(4\pi |c_0|)^{\ell_1'} e^{\ell_1'} \ell_1'!(\ell_1'+1)\left(1+\frac{1+e^{-2\pi  \Real(q_0)}}{\varepsilon}\right)^{\ell_1'}\sum_{|\ell|=\gamma}\binom{\gamma}{\ell}2\pi e^{2\pi|\Real(q)|}(4\pi C_{b,a_0})^{\ell_2'}\\
    & \phantom{\leq} \times |\xi|^{\ell_1}e^{-\frac{L}{2}|\xi|^{\frac{1}{\mu}}}C_{1,b}^{\ell_1}\ell_1!^\sigma (1+2C_{0,b})^{\ell_1}  \sup_{(t,\xi)\in\g}e^{L|\xi|^{\frac{1}{\mu}}}|\partial_\xi^{\ell'_3}\partial_t^{\ell_2}\widehat{f}(t,\xi)|\\
    &\leq \varepsilon^{-1} \sum_{|\ell'|=\beta} \sum_{|\ell|=\gamma}\binom{\beta}{\ell'}\binom{\gamma}{\ell}K_{c_0,q_0}^{\ell_1'}
    (4\pi C_{b,a_0})^{\ell_2'}C_{1,b}^{\ell_1} (1+2C_{0,b})^{\ell_1}\textstyle{(\frac{2\mu}{L})^{\mu\ell_1}}\ell_1'!\ell_1!^\mu\ell_1!^\sigma C_{0,f}C_{1,f}^{\ell_3'+\ell_2}\ell_2!^{\sigma}\ell_3'!^\mu\\
    &\leq C_0 \tilde{C}_{1}^{\beta+\gamma} 3^\beta 2^\gamma \beta!^{\mu}\gamma!^{\sigma+\mu}\leq C_0C_1^{\beta+\gamma}\beta!^\mu\gamma!^{\sigma+\mu},
\end{align*}
for every $\xi<0$, where $C_0=\varepsilon^{-1}C_{0,f}$, $\tilde{C}_1=K_{c_0,q_0}\max\{4\pi C_{b,a_0},1\}(1+2C_{0,b})(\frac{2\mu}{L})^\mu C_{1,b}C_{1,f},$ $C_1=6\tilde{C}_1$, and $K_{c_0,q_0}= \max\{4\pi|c_0|e\left(1+\varepsilon^{-1}\left(1+e^{-2\pi  \Real(q_0))}\right)\right),1\}$, .

For $\xi\geq0$, we obtain similar inequalities for $\widehat{u}(t,\xi)$ by an analogous argument. 

This proves that $\widehat{u}\in \mathcal{S}_{\sigma+\mu,\mu}(\g)$.  Consequently $u\in\mathcal{S}_{\sigma+\mu,\mu}\subset  \Smu$ by Corollary \ref{coro_gelf_transf} and therefore $P$ is $\mathscr{F}_{\mu}$-globally hypoelliptic. In the case where $b(t)\leq 0$ and so $b_0<0$, the proof is similar, only now we swap the arguments for $\widehat{u}(t,\xi)$ with respect to the cases $\xi\geq 0$ and $\xi<0$.
\end{proof}

\end{document}
